\begin{thebibliography}{20}
\bibitem{ref1}Eric Allender. The new complexity landscape around circuit minimization. In Alberto Leporati, Carlos Martín-Vide, Dana Shapira, and Claudio Zandron, editors, Language and Automata Theory and Applications - 14th International Conference, LATA 2020, Milan, Italy, March 4-6, 2020, Proceedings, volume 12038 of Lecture Notes in Computer Science, pages 3–16. Springer, 2020. doi:10.1007/978-3-030-40608-0\_1.
\bibitem{ref2}Eric Allender, Harry Buhrman, Michal Koucký, Dieter van Melkebeek, and Detlef Ronneburger. Power from random strings. SIAM J. Comput., 35(6):1467–1493, 2006.
\bibitem{ref3}Eric Allender and Bireswar Das. Zero knowledge and circuit minimization. Inf. Comput., 256:2–8, 2017.
\bibitem{ref4}Eric Allender, Michal Koucký, Detlef Ronneburger, and Sambuddha Roy. The pervasive reach of resource-bounded kolmogorov complexity in computational complexity theory. Journal of Computer and System Sciences, 77(1):14–40, 2011.
\bibitem{ref5}David Buchfuhrer and Christopher Umans. The complexity of boolean formula minimization. J. Comput. Syst. Sci., 77(1):142–153, 2011.
\bibitem{ref6}Marco L. Carmosino, Russell Impagliazzo, Valentine Kabanets, and Antonina Kolokolova. Learning algorithms from natural proofs. In Proceedings of the 31st Conference on Computa- tional Complexity, 2016.
\bibitem{ref7}Alexander Golovnev, Rahul Ilango, Russell Impagliazzo, Valentine Kabanets, Antonina Koloko- lova, and Avishay Tal. Ac0 [p] lower bounds against MCSP via the coin problem. In ICALP, volume 132 of LIPICS, pages 66:1–66:15, 2019.
\bibitem{ref8}Shuichi Hirahara. Non-black-box worst-case to average-case reductions within NP. In 59th IEEE Annual Symposium on Foundations of Computer Science, FOCS, pages 247–258, 2018.
\bibitem{ref9}Shuichi Hirahara, Igor Carboni Oliveira, and Rahul Santhanam. NP-hardness of minimum circuit size problem for OR-AND-MOD circuits. In 33rd Computational Complexity Conference, CCC, volume 102, pages 5:1–5:31, 2018.
\bibitem{ref10}Valentine Kabanets and Jin-Yi Cai. Circuit minimization problem. In Proceedings of the Thirty-Second Annual ACM Symposium on Theory of Computing, STOC ’00, page 73–79, 2000.
\bibitem{ref11}S. A. Lozhkin. Tighter bounds on the complexity of control systems from some classes. Mat. Voprosy Kibernetiki 6, pages 189–214 (in Russian), 1996.
\bibitem{ref12}Oleg B. Lupanov. Complexity of formula realization of functions of logical algebra. Problemy Kibernetiki, 3:61–80, 1960.
\bibitem{ref13}William J. Masek. Some NP-complete set covering problems. Unpublished Manuscript, 1979.
\bibitem{ref14}Cody D. Murray and R. Ryan Williams. On the (non) np-hardness of computing circuit complexity. Theory of Computing, 13(1):1–22, 2017.
\bibitem{ref15}Erkki Mäkinen. Generating random binary trees — a survey. Information Sciences, 115(1):123– 136, 1999.
\bibitem{ref16}Nicholas Pippenger. Information theory and the complexity of boolean functions. Mathematical Systems Theory, 10:129–167, January 1977.
\bibitem{ref17}Alexander A. Razborov and Steven Rudich. Natural proofs. J. Comput. Syst. Sci., 55(1):24–35, 1997.
\bibitem{ref18}Michael Rudow. Discrete logarithm and minimum circuit size. Inf. Process. Lett., 128:1–4, 2017.
\bibitem{ref19}Rahul Santhanam. Pseudorandomness and the minimum circuit size problem. In Thomas Vidick, editor, 11th Innovations in Theoretical Computer Science Conference, ITCS, volume 151 of LIPIcs, pages 68:1–68:26, 2020.
\bibitem{ref20}B. A. Trakhtenbrot. A survey of russian approaches to perebor (brute-force searches) algorithms. IEEE Ann. Hist. Comput., 6(4):384–400, October 1984.
\end{thebibliography}
